\newif\ifglobal

%%%%%%%%%%%%%%%%%%%%%%%
%%% Compile to view %%%
%%%%%%%%%%%%%%%%%%%%%%%

%%%%%%%%%%%%%%%%%%%%%%%%%%%%%%%%%%%%%%%%%%%%%%%%%%%
%%% Readme:                                     %%%
%%% https://www.overleaf.com/read/bwdjvfjksvjy/ %%%
%%%%%%%%%%%%%%%%%%%%%%%%%%%%%%%%%%%%%%%%%%%%%%%%%%%

% >>> M?: marks a module setting number or important notice

% >>> M4: Set {./relative-path} to external files for this module <<<
\def \modulefiles {./27-RNG}
% To add an external file, use:
% (Replace '---' with actual filename)
% '\input{\modulefiles/tables/---.tex}' for tables
% '\includegraphics{\modulefiles/figures/---}' for figures
% '\subfile{\modulefiles/---}' for register files

% >>> M1: This module is in global project? <<<
% 'YES' - uncomment; 'NO' - comment out
\globaltrue

\ifglobal
    \documentclass[main\_\_CN.tex]{subfiles}

\else
    % Font "FandolSong-Regular" used by ctex does not have GSUB table.
    % It causes warnings that do not affect compilation and can be ignored.

    \PassOptionsToPackage{quiet}{fontspec}

    \documentclass[openany, 10pt]{book}

    % >>> M2: In ./00-shared/preamble-trm-module, also find
    %         settings for visibility of todo notes, watermark <<<
    \usepackage{preamble-trm-module}


    % Enable Chinese typesetting
    \usepackage{ctex}

    % >>> M3: Temporary labels in Chinese
    %         you can add your own labels there <<<
    \myexternaldocument{./00-shared/temp-labels-cn}

\fi %\ifglobal


\setcounter{secnumdepth}{4}




% >>> M5: If this module needs any updates in
% './00-shared/preamble-trm-module.sty'
% (include additional package, etc.), contact your module PM
% or create the file './00-shared/changelog.tex' make the
% necessary updates yourself and document them in this file <<<

% Make text left-aligned, comment out for full justification
\raggedright


\begin{document}


% Sets language into which pre-defined bits of text are translated
\selectlanguage{chinese}

% Do not delete this block
\ifglobal
% Add the following front matter if
% this module is outside of global project
\else
    \InputIfFileExists{readme}{}{}
    {\let\clearpage\relax \listoftodos}
    {\let\clearpage\relax \listofdonetodos}
    \tableofcontents
    %\listoftables
    %\listoffigures
\fi


%%%%%%%%%%%%%%%%%%%%%%%%%
%%% TRM Module Begins %%%
%%%%%%%%%%%%%%%%%%%%%%%%%

\hypertarget{rng}{}
\chapter{随机数发生器 (RNG)}\label{mod:rng}

\section{概述}

\chipname{} 内置一个真随机数发生器，其生成的 32 位随机数可作为加密等操作的基础。

\section{主要特性}

\chipname{} 的随机数发生器可通过物理过程而非算法生成真随机数，所有生成的随机数在特定范围内出现的概率完全一样。

\section{功能描述}\label{sec:rng-func-descr}

系统可以从随机数发生器的寄存器 \hyperref[regdesc:LPPERIRNGDATAREG]{LPPERI\_RNG\_DATA\_REG} 中读取随机数，每个读到的 32 位随机数都是真随机数，噪声源为系统中的\textbf{热噪声}和\textbf{异步时钟}。

\begin{itemize}
    \item \textbf{热噪声}可以来自 SAR ADC 或高速 ADC 或两者兼有。当芯片的 SAR ADC 或高速 ADC 工作时，就会产生比特流，并通过异或 (XOR) 逻辑运算作为随机数种子进入随机数发生器。
    \item RC\_FAST\_CLK 是一种\textbf{异步时钟源}，会产生电路亚稳态。这种亚稳态可以增加随机数发生器的熵值。
\end{itemize}


\begin{figure}[H]
    \centering
    \includegraphics[width=0.9\textwidth]{\modulefiles/figures/29-RNG-noise-source.pdf}
    \caption{噪声源}
    \label{fig:rng-noise-source}
\end{figure}


当高速 ADC 打开时，随机数发生器在每个 APB 时钟周期（通常为 32 MHz）内获得 2 位的熵。

当 SAR ADC 打开时，随机数发生器在每个 SAR ADC 采样周期内获得 2 位熵值。SAR ADC 采样周期通常为 66 KHz，由 SAR ADC 的配置决定。

我们在仅打开高速 ADC 的状态下，以 5MHz 的速率从\hyperref[regdesc:LPPERIRNGDATAREG]{LPPERI\_RNG\_DATA\_REG} 读取了 2GB 的数据样本，并使用 Dieharder 随机数测试套件（版本 3.31.1）对样本进行了测试。最终，样本通过了所有测试。

\section{编程指南}

在使用 \chipname{} 的随机数发生器时，应该至少保证 SAR ADC 或 高速 ADC\hyperref[others:rng-note-clock-seed]{$^1$} 处于使能状态,同时建议使能 RC\_FAST\_CLK\hyperref[others:rng-note-clock-seed]{$^2$} 以增加进入随机数的熵值，否则可能会导致产生伪随机数，应注意避免。其中，
\begin{itemize}
    \item SAR ADC 受控于 DIG ADC 控制器。详见 \ref{mod:adc} \textit{\nameref{mod:adc}} 章节。

    建议 SAR ADC 的配置如下：
    \begin{itemize}
        \item 设置采集芯片内部电压的采样通道
         \item 设置控制器使用定时器驱动的连续采样模式
    \end{itemize}
    \item 高速 ADC 在无线模块开启时自动打开。
    \item RC\_FAST\_CLK 默认使能，无法关闭。


\end{itemize}
\begin{tiplisting}
\vspace{-2.5em}
\begin{enumerate}
    \label{others:rng-note-clock-seed}
    \item 注意，在无线模块开启时，极端情况下高速 ADC 有读值饱和的可能，这会降低熵值。因此，建议在 无线模块开启时，同时打开 SAR ADC 产生随机数。
    \item RC\_FAST\_CLK 时钟仅可以提高随机数发生器的墒值。然而，为了保证随机数发生器可以获得足够大的熵值，仍建议在使用随机数发生器时至少保证 SAR ADC 或高速 ADC 处于工作状态。
\end{enumerate}
\end{tiplisting}

在使用随机数发生器时，请多次读取 \hyperref[regdesc:LPPERIRNGDATAREG]{LPPERI\_RNG\_DATA\_REG} 寄存器的值，直至获得足够多的随机数。
请确保读取该寄存器的频率不超过上文 \ref{sec:rng-func-descr} 小节中描述的熵生成频率。

%%%%%%%%%%%%%%%%%%%%%%%%
%%% Register Summary %%%
%%%%%%%%%%%%%%%%%%%%%%%%

\hypertarget{rng-reg-summ}{}
\section{寄存器列表}
\label{sec:rng-reg-summ}
请查看章节 \hyperref[glossary-access-types]{寄存器的访问类型}，了解“访问”列缩写的含义。


\begin{table}[H]
    \centering
    \begin{tabular}{ | l | l | l | l |  }
    \hline
    \rowcolor{lightgray}
    \textbf{名称} & \textbf{描述} & \textbf{地址} & \textbf{访问} \\
    \hline
    \hyperref[regdesc:LPPERIRNGDATAREG]{LPPERI\_RNG\_DATA\_REG} & 随机数数据 & 0x600B\_280C & 只读 \\
    \hline
    \end{tabular}
\end{table}



%%%%%%%%%%%%%%%%%
%%% Registers %%%
%%%%%%%%%%%%%%%%%


\section{寄存器}

\begin{register}{H}{LPPERI\_RNG\_DATA\_REG}{0x600B\_280C}\label{regdesc:LPPERIRNGDATAREG}
\regfield{LPPERI\_RNG\_DATA}{32}{0}{{0x00000000}}%
\reglabel{Reset}\regnewline%
\begin{regdesc}\begin{reglist}
\item [LPPERI\_RNG\_DATA]随机数来源。（只读）
\end{reglist}\end{regdesc}
\end{register}


\end{document}

